\subsection{Etap Deploy}
Faza wdrożenia (\textit{Deploy}) odpowiada za przeniesienie artefaktów na środowisko uruchomieniowe oraz zapewnienie odpowiedniej infrastruktury.

\subsubsection{Środowisko uruchomieniowe: Kubernetes (k3d)}
Głównym środowiskiem uruchomieniowym mikroserwisów jest klaster \textbf{Kubernetes} (K8s) — otwartoźródłowy system automatyzacji wdrażania, skalowania i zarządzania skonteneryzowanymi aplikacjami. Pełni rolę orkiestratora: dba o to, aby zdefiniowana liczba instancji aplikacji (podów) była zawsze uruchomiona i dostępna, a w przypadku awarii automatycznie je restartuje. Lokalne klastry Kubernetes oferują niemal pełną kompatybilność z produkcyjnymi klastrami chmurowymi (EKS, GKE), co pozwala testować manifesty, polityki sieciowe i mechanizmy skalowania bezpośrednio na maszynie programisty.

\begin{techcard}{Kubernetes}

\textbf{Opis:} Kubernetes (K8s) to otwartoźródłowy system automatyzacji wdrażania, skalowania i zarządzania skonteneryzowanymi aplikacjami. Pełni rolę orkiestratora, dbając o to, aby zdefiniowana liczba instancji aplikacji (podów) była zawsze uruchomiona i dostępna dla użytkowników.

\textbf{Dobre nawyki:}
\begin{itemize}
  \item \textbf{Zarządzanie zasobami:} Obowiązkowe definiowanie \texttt{requests} (zasoby gwarantowane) i \texttt{limits} (maksymalne zużycie) zapobiega wzajemnemu zagłuszaniu się usług (efekt \textit{noisy neighbor}).
  \item \textbf{Sondy stanu podów:} \texttt{livenessProbe} (czy kontener żyje) oraz \texttt{readinessProbe} (czy może przyjmować ruch) są kluczowe dla samonaprawialności klastra i bezprzestojowych wdrożeń (\textit{Rolling Updates}).
  \item \textbf{Podejście deklaratywne (GitOps):} Konfiguracja przechowywana w repozytorium i nakładana przez \texttt{kubectl apply}, a nie komendy imperatywne.
  \item \textbf{Network Policies:} Domyślne ograniczanie komunikacji między podami do niezbędnego minimum -- zasada \textit{Zero Trust}.
\end{itemize}

\textbf{Rola w projekcie:} Kubernetes stanowi docelowe środowisko uruchomieniowe. O ile Docker pozwala stworzyć obraz, o tyle Kubernetes odpowiada za to, by ten obraz funkcjonował w klastrze, był dostępny pod konkretnym adresem i automatycznie restartował się w przypadku awarii.

\end{techcard}

\begin{lstlisting}[language=yaml, caption={Manifest Kubernetes Deployment z sondami kondycji i limitami zasobów (infrastructure/k8s/users-deployment.yaml)}, label={lst:k8s-deployment}]
apiVersion: apps/v1
kind: Deployment
metadata:
  name: users-deployment
spec:
  replicas: 3
  selector:
    matchLabels:
      app: users
  template:
    metadata:
      labels:
        app: users
      annotations:
        prometheus.io/scrape: "true"
        prometheus.io/path: /q/metrics
        prometheus.io/port: "8080"
    spec:
      containers:
      - name: users
        image: mateuszwas/users:1.0.0
        ports:
        - containerPort: 8080
        resources:
          requests:
            memory: "64Mi"
            cpu: "50m"
          limits:
            memory: "256Mi"
            cpu: "500m"
        livenessProbe:
          httpGet:
            path: /q/health/live
            port: 8080
        readinessProbe:
          httpGet:
            path: /q/health/ready
            port: 8080
\end{lstlisting}

W części praktycznej lokalny klaster uruchamia narzędzie \textbf{k3d} — cienka warstwa opakowująca \textbf{k3s} (certyfikowaną przez CNCF dystrybucję Kubernetes) działającą wewnątrz kontenerów Docker. Kluczowym wymaganiem było dotarcie do Ingressu bezpośrednio z hosta bez dodatkowych komend. Alternatywa — \textbf{minikube} z driverem Docker — wymaga uruchomienia \texttt{minikube tunnel} w osobnym terminalu, by ruch z hosta docierał do zasobów Ingress. k3d rozwiązuje ten problem przez mapowanie portów już na etapie tworzenia klastra:

\begin{lstlisting}[language=bash, caption={Tworzenie klastra k3d z mapowaniem portów Ingress (infrastructure/k3d/cluster-create.sh)}, label={lst:k3d-create}]
k3d cluster create bsc-devops \
  --agents 2 \
  --port "80:80@loadbalancer" \
  --port "443:443@loadbalancer"
\end{lstlisting}

Powyższe polecenie mapuje port 80 hosta na load balancer k3d, dzięki czemu każdy zasób Ingress jest natychmiast dostępny pod adresem \texttt{http://localhost}. Klaster startuje w czasie poniżej 30 sekund (w porównaniu do 1--2 minut dla minikube).

\begin{itemize}
    \item \textbf{Wbudowany kontroler Ingress:} k3s domyślnie instaluje \textbf{Traefik} jako kontroler Ingress — nie wymaga ręcznego instalowania dodatków.
    \item \textbf{Wiele węzłów:} Flaga \texttt{--agents N} pozwala tworzyć klastry wielowęzłowe, umożliwiając testowanie zachowania schedulera zbliżonego do środowisk produkcyjnych.
    \item \textbf{Import obrazów lokalnych:} Polecenie \texttt{k3d image import <obraz> -c bsc-devops} ładuje lokalne obrazy Docker bezpośrednio do klastra bez potrzeby wypychania ich do zewnętrznego rejestru.
\end{itemize}

\begin{techcard}{k3d}

\textbf{Opis:} k3d to narzędzie uruchamiające klastry \textbf{k3s} (lekkiej dystrybucji Kubernetes certyfikowanej przez CNCF) wewnątrz kontenerów Docker. Każdy węzeł klastra jest osobnym kontenerem, a porty load balancera są mapowane na host już w momencie tworzenia klastra. Rozwiązuje to główną bolączkę alternatywnych narzędzi lokalnych -- konieczność uruchamiania dodatkowych procesów tunelujących ruch Ingress.

\textbf{Dobre nawyki:}
\begin{itemize}
  \item \textbf{Mapowanie portów przy tworzeniu:} Flagi \texttt{--port} należy podać w poleceniu \texttt{k3d cluster create} -- nie da się ich dodać do działającego klastra bez jego rekonstrukcji.
  \item \textbf{Import obrazów lokalnych:} Zamiast wypychać obraz do zewnętrznego rejestru podczas iteracji deweloperskiej, należy użyć \texttt{k3d image import <obraz> -c <klaster>} -- ładuje obraz bezpośrednio do containerd w węzłach klastra.
  \item \textbf{Domyślny Ingress: Traefik:} k3s dostarcza Traefik zamiast nginx; manifesty Ingress muszą używać \texttt{ingressClassName: traefik}.
  \item \textbf{Wielowęzłowość:} Dla testów zbliżonych do produkcji warto dodać agentów (\texttt{--agents 2}), by sprawdzić zachowanie schedulera przy ograniczeniach \texttt{affinity}/\texttt{tolerations}.
\end{itemize}

\textbf{Rola w projekcie:} k3d stanowi lokalne środowisko deweloperskie i demonstracyjne dla całego stosu mikroserwisów. Zastąpił minikube ze względu na natywne mapowanie portów Ingress (brak \texttt{minikube tunnel}), czas startu klastra poniżej 30 sekund oraz pełną skryptowalność tworzenia klastra jedną komendą CLI.

\end{techcard}

\subsubsection{Automatyzacja wdrożeń: GitOps (ArgoCD)}
Zamiast tradycyjnych narzędzi \textit{Configuration Management} (jak Ansible), zdecydowano się na podejście \textbf{GitOps} przy użyciu narzędzia \textbf{ArgoCD} (ADR-0013). GitOps zakłada, że repozytorium Git jest jedynym źródłem prawdy o stanie infrastruktury -- kontroler stale monitoruje zmiany w chartach Helm i automatycznie synchronizuje stan klastra z definicją w kodzie, eliminując potrzebę ręcznego wykonywania poleceń \texttt{kubectl} z poziomu CI/CD.

\begin{techcard}{ArgoCD}

\textbf{Opis:} ArgoCD to deklaratywny kontroler GitOps dla Kubernetes, będący projektem CNCF o statusie \textit{graduated}. Działa w modelu \textbf{pull-based}: nasłuchuje zmian w repozytorium Git lub rejestrze OCI i automatycznie synchronizuje stan klastra do zdefiniowanego tam stanu docelowego. Centralną abstrakcją jest zasób \texttt{Application}, który wiąże źródło (chart Helm, katalog manifestów) z docelowym klastrem i przestrzenią nazw. Wbudowany interfejs webowy pozwala w czasie rzeczywistym obserwować drzewo zasobów aplikacji oraz historię wdrożeń.

\textbf{Dobre nawyki:}
\begin{itemize}
  \item \textbf{Automated sync z \texttt{selfHeal}:} Włączenie \texttt{selfHeal: true} gwarantuje, że każda ręczna zmiana w klastrze zostanie natychmiast cofnięta -- klaster jest zawsze zgodny z tym, co opisuje Git.
  \item \textbf{\texttt{prune: true} z rozwagą:} Automatyczne usuwanie zasobów odłączonych od charta skraca dług konfiguracyjny, ale wymaga pewności, że wszystkie zasoby są zarządzane przez ArgoCD.
  \item \textbf{Application CRD w Git:} Plik \texttt{gitops/application.yaml} powinien być przechowywany w repozytorium obok charta -- cały opis dostawy jest wówczas wersjonowany i audytowalny.
  \item \textbf{Pinowanie \texttt{targetRevision}:} Używanie konkretnego numeru wersji charta (nie \texttt{latest}) zapewnia powtarzalność wdrożeń i pełną historię zmian.
\end{itemize}

\textbf{Rola w projekcie:} ArgoCD zamyka pętlę CI/CD: GitHub Actions buduje obraz, publikuje chart do GHCR, a ArgoCD automatycznie wykrywa nową wersję charta i synchronizuje klaster k3d. Zasób \texttt{Application} wskazuje na chart \texttt{bsc-devops} w GHCR i celuje w lokalny klaster k3d. Wbudowany interfejs webowy dostarcza wizualizację drzewa zasobów przydatną jako materiał ilustracyjny pracy.

\end{techcard}

\begin{lstlisting}[language=yaml, caption={ArgoCD Application CRD -- deklaracja dostawy charta \texttt{bsc-devops} z GHCR do klastra k3d (gitops/application.yaml)}, label={lst:argocd-application}]
apiVersion: argoproj.io/v1alpha1
kind: Application
metadata:
  name: bsc-devops
  namespace: argocd
spec:
  project: default
  source:
    repoURL: ghcr.io/mateuszwas/charts   # OCI Helm registry (GHCR)
    chart: bsc-devops
    targetRevision: "0.5.0"
    helm:
      valueFiles:
        - values.yaml
  destination:
    server: https://kubernetes.default.svc
    namespace: bsc-devops
  syncPolicy:
    automated:
      prune: true      # usuwa zasoby usunięte z charta
      selfHeal: true   # cofa ręczne zmiany niezgodne z Gitem
    syncOptions:
      - CreateNamespace=true
\end{lstlisting}
